\chapter{Introduction}

The legal industry is increasingly adopting large language models (LLMs) to assist lawyers in navigating the complexity of modern litigation.
Legal cases are inherently dynamic: they evolve over time as new documents, correspondence, evidence, and procedural developments continuously
emerge. What was once a central claim may become peripheral; a newly submitted document may reframe the entire dispute. This progression
means that not only does the volume of material grow, but the \textit{relevance and weight} of existing information shifts as the case develops. \\ \\

A fundamental limitation of current LLM-based systems is that they are stateless. Without explicit mechanisms for maintaining case state,
models must re-read the full document history for every query. More critically, even models that handle large contexts tend to perform well
in early stages but \textit{drift} over extended, multi-turn interactions --- gradually losing track of what is current, what has been superseded,
and what actually matters. In legal case management, this drift is not merely an inconvenience; it can lead to factually incorrect summaries,
missed deadlines, or misrepresentation of parties and claims. \\ \\

This thesis addresses this problem through two complementary mechanisms: \textit{compression} of case-relevant information into a structured
representation, and \textit{state maintenance} --- the continuous, structured tracking of case essentials as the project evolves. The vehicle
for this is a structured \textit{FactSheet}: a living, machine-readable summary of the case that is initialized from initial input and updated
incrementally as new documents and information arrive. \\ \\

The research investigates whether maintaining an explicit, structured state representation improves an LLM agent's ability to preserve
contextual accuracy over the full lifecycle of a legal case --- particularly as conversation length grows and the document landscape changes.
The system is implemented as a LangGraph-based conversational agent, incorporating multi-phase document processing pipelines for both
project initialization and incremental updates. \\ \\

The contributions of this thesis are: (1) the design and implementation of initialization and update pipelines for structured legal case
state management, (2) an evaluation framework comparing structured state agents against RAG-only and baseline approaches, and (3) an
analysis of accuracy degradation, compression efficiency, and state consistency across extended multi-turn conversations. \\ \\

The remaining chapters are organized as follows: Chapter 2 presents background and related work on long-context LLMs, context compression,
and retrieval-augmented generation. Chapter 3 describes the methodology, system architecture, and evaluation design. Chapter 4 presents
experimental results. Chapter 5 discusses findings and implications, and Chapter 6 concludes with a summary and directions for future work.